\documentclass{rmf-d}
\usepackage{nopageno,multicol,times,epsf,amsmath,amssymb}
\usepackage[T1]{fontenc} 
\usepackage[spanish, mexico]{babel}
\usepackage[]{caption}
\usepackage{graphicx}
\usepackage{blindtext}
\usepackage{color}
\usepackage[hidelinks]{hyperref}
\usepackage[backend=biber, natbib=true,style=numeric]{biblatex}
\usepackage{hhline}
\usepackage{xcolor}
\usepackage{subcaption}
\usepackage{cancel}
\usepackage{afterpage}
\usepackage{amsthm}
\usepackage{bookmark}
\usepackage{multirow}
\usepackage{makecell}
\usepackage{booktabs}
\usepackage{pgfplots}
\usepackage{circuitikz}
\usepackage{physics}
\usepackage{pgfplots}
\usepackage{listings}





\definecolor{codegreen}{rgb}{0,0.6,0}
\definecolor{codegray}{rgb}{0.5,0.5,0.5}
\definecolor{codepurple}{rgb}{0.58,0,0.82}
\definecolor{backcolour}{rgb}{0.95,0.95,1}
\definecolor{airforceblue}{rgb}{0.36, 0.54, 0.66}
\definecolor{americanrose}{rgb}{1.0, 0.01, 0.24}
\definecolor{antiquefuchsia}{rgb}{0.57, 0.36, 0.51}
\definecolor{amaranth}{rgb}{0.9, 0.17, 0.31}
\definecolor{brightmaroon}{rgb}{0.76, 0.13, 0.28}
\definecolor{blue(pigment)}{rgb}{0.2, 0.2, 0.6}
\lstdefinestyle{mystyle}{
  backgroundcolor=\color{backcolour},   commentstyle=\color{codegreen},
  keywordstyle=\color{magenta},
  numberstyle=\tiny\color{codegray},
  stringstyle=\color{codepurple},
  basicstyle=\ttfamily\footnotesize,
  breakatwhitespace=false,         
  breaklines=true,                 
  captionpos=b,                    
  keepspaces=true,                 
  numbers=left,                    
  numbersep=5pt,                  
  showspaces=false,                
  showstringspaces=false,
  showtabs=false,                  
  tabsize=2
}

\lstset{style=mystyle}





\renewcommand*\descriptionlabel[1]{\hspace\leftmargin$#1$}

\addbibresource{ref.bib}
\bibliography{ref}


\pgfplotsset{compat=1.17}
\renewcommand\theadalign{cc}
\renewcommand\theadfont{}
\renewcommand\theadgape{\Gape[4pt]}
\renewcommand\cellgape{\Gape[0pt]}
\spanishdecimal{.}
\def\rmfcornisa{Métodos numéricos \hfill Cuarto Semestre\
\hfill Marzo 2026}
\newcommand{\ssc}{\scriptscriptstyle}
\definecolor{LightGray}{gray}{0.9}


\def\rmfcintilla{{\it Facultad de Ciencias Físico Matemáticas\/}}
\clearpage \pagestyle{myheadings}
\setcounter{page}{1}


















\begin{document}
%\markboth{Universidad Autónoma de Coahuila}




\title{Método de Secante
\vspace{-6pt}}
\author{Javier Alexander López Contreras }

\address{Facultad de Ciencias Físico Matemáticas (FCFM), Universidad Autónoma de Coahuila, Campus Campo redondo }
\\

\maketitle




\begin{center}
\recibido{2 de marzo de 2026
\vspace{-10pt}}
\end{center}


\begin{abstract}
En este documento se analizará el código en Phyton del método de Secante; que es un método que se utiliza para encontrar la raíz de una función . Se busca explicar de mejor manera el código.
\end{abstract}











\begin{multicols}{2}

\section{INTRODUCCIÓN}
\label{sec:1}
El estudio de los métodos numéricos es esencial para hallar soluciones a ecuaciones trascendentes y polinomiales donde los métodos analíticos resultan insuficientes. El enfoque principal de este trabajo es el Método de la Secante, un algoritmo de búsqueda de raíces que destaca por su eficiencia al no requerir el cálculo de la derivada de la función, a diferencia de otros métodos abiertos como Newton-Raphson. El objetivo primordial es analizar su comportamiento, velocidad de convergencia y estabilidad ante diferentes tipos de funciones matemáticas, utilizando un soporte computacional para observar su evolución iterativa.
































\section{ALGORITMO Y EXPLICACIÓN DEL
MÉTODO DE SECANTE}
\label{sec:2}
El método de la secante se basa en una aproximación lineal de la función. En lugar de utilizar la recta tangente, se utiliza una \textbf{recta secante} que pasa por dos puntos iniciales $x_0$ y $x_1$.

El proceso lógico sigue estos pasos:
\begin{itemize}
    \item \textbf{Interpolación lineal}: Se traza una línea recta entre los punto $(x_0,f(x_0))$ y $(x_1,f(x_1))$
    \item \textbf{Intersección con el eje $x$}: El punto donde esta recta cruza el eje horizontal se define como la nueva aproximación $x_2$
    \item \textbf{Criterio de recurrencia}: La fórmula matemática que hace todo este avance es:
    \begin{equation}
        x_i+1=x_1-\frac{f(x_i)(x_i-x_{i-1})}{f(x_i)-f(x_{i-1})}
    \end{equation}
    \item \textbf{Condición de parada}: El proceso se repite hasta que la diferencia entre dos aproximaciones sucesivas sea menor a la tolerancia predefinida, o se alcance un número máximo de iteraciones. 
\end{itemize}






























\section{ANÁLISIS DEL CÓDIGO}
\label{sec:3}
Para validar el método, se desarrollo un programa que traduce la lógica matemática a instrucciones de control. Los puntos clave son:
\begin{itemize}
    \item \textbf{Evaluación dinámica de la función}

    Para que el método de la secante sea universal, el programa debe interpretar la entrada del usuario como una operación matemática real.
    
    \begin{lstlisting}[language=Python, caption=interpretación de la funcion]
    def f(x):
    return eval(funcion.get())
    \end{lstlisting}

    \item \textbf{Implementación de la fórmula de la secante.}

    Dentro del bucle while, se traduce la ecuación de recurrencia para hallar el siguiente punto de aproximación.
    \begin{lstlisting}[language=Python, caption=formula de la secante]
x2 = x1 - (fx1 * (x1 - x0)) / (fx1 - fx0)    
    \end{lstlisting}
    \item \textbf{Manejo crítico de la pendiente.}

    El método de la secante falla si la pendiente entre los dos puntos es nula. Se implementó este control preventivo.
    \begin{lstlisting}[language=Python, caption=manejo de division entre 0]
    if abs(fx1 - fx0) < 1e-15:
    resultado.config(text='Error: División por cero (pendiente horizontal)', fg='red')
    return
    \end{lstlisting}

    \item \textbf{Actualización de Estados de Iteración.}

    Para que el algoritmo avance hacia la raíz, los valores deben desplazarse en cada ciclo.
    \begin{lstlisting}[language=Python, caption=Actualización de x]
     x0 = x1
    x1 = x2
    iter_count += 1
    \end{lstlisting}
\end{itemize}




















\section{RESULTADOS Y ANÁLISIS}
\label{sec:4}
Se seleccionaron tres tipos de pruebas para evaluar los límites del método de la secante.

\subsection{Convergencia rápida}
Se seleccionó la función $x^2-4$ para verificar la exactitud en polinomios y validar que el algoritmo detecta raíces enteras sin oscilaciones.
\begin{figure}[H]
    \centering
    \includegraphics[width=0.5\linewidth]{prueba_sec_1.png}
    \caption{Prueba realizada para la convergencia rápida}
    \label{fig:placeholder}
\end{figure}

\subsection{Trascendente}
Se usó la función $e^{-x} - x$  para evaluar funciones exponenciales. Esta prueba demuestra la ventaja sobre el método de bisección, esto por la velocidad de convergencia y aprovechamiento de la magnitud de $f(x)$

\begin{figure}[H]
    \centering
    \includegraphics[width=0.5\linewidth]{prueba_sec_2.png}
    \caption{Prueba usando e}
    \label{fig:placeholder}
\end{figure}

\subsection{Error de pendiente.}
Para esta prueba, se usó la función $cos(x)$, esto da una división entre 0, ya que estamos forzando una pendiente horizontal teniendo $x_0=-1$ y $x_1=1$, donde la función no toca el eje x

\begin{figure}[H]
    \centering
    \includegraphics[width=0.5\linewidth]{prueba_sec_3.png}
    \caption{Forzando una pendiente horizontal}
    \label{fig:placeholder}
\end{figure}































\section{CONCLUSIONES}
\label{sec:5}
El análisis permite concluir que el método de la secante es una de las herramientas más potentes en el análisis numérico debido a su rapidez y su simplicidad. A través de la implementación realizada, se demostró que el éxito del método depende de la continuidad de la función y de la selección de los puntos iniciales. 



























\end{multicols}
\medline
\begin{multicols}{2}

    
\printbibliography


\end{multicols}
\end{document}
